\chapter{この本の使い方と記号}
\label{ch:tb-guide}

\section{何を収めたか}

専門基礎科目は、収録した5年度すべてで
\textbf{線形代数}・\textbf{微分積分}・\textbf{位相／距離空間}の三分野から出題されている。
本書はその三分野を、次の順に積み上げる。

\begin{enumerate}
  \item 第\ref{ch:tb-foundations}章：論理・集合・写像・実数・数列。三分野すべての土台。
  \item 第\ref{ch:tb-functions}章：初等関数と暗記すべき公式。$\sec$ や双曲線関数の約束もここ。
  \item 第\ref{ch:tb-differential}章・第\ref{ch:tb-integral}章：一変数の微分と積分。
  \item 第\ref{ch:tb-multivariable}章：多変数の微分積分。重積分、変数変換、Gamma・Beta 関数。
  \item 第\ref{ch:tb-linear-one}章・第\ref{ch:tb-linear-two}章：線形代数。空間と行列式、固有値と内積。
  \item 第\ref{ch:tb-topology}章：距離空間と位相空間。コンパクト性と連結性。
\end{enumerate}

各節の終わりには \textbf{Memo} を置き、その道具が過去問のどこで実際に使われたかを示した。
問題番号は解答編の表記（例：2021年度 第1期 専門基礎 第2問）に合わせている。

\section{読み方の約束}

\begin{itemize}
  \item \textbf{Def.}（定義）は言葉の意味を定める。定義は証明の対象ではなく、覚えて使うものである。
  \item \textbf{Thm.}（定理）・\textbf{Prop.}（命題）は証明できる主張である。
        入試で「示せ」と問われうるものには証明または導出の筋を付けた。
  \item \textbf{Ex.}（例）は、定義や定理が具体的にどう動くかを見る場所である。
  \item \textbf{Rem.}（注意）は、間違えやすい点をまとめた。
\end{itemize}

証明を省いた定理は、答案で「$\sim$ の定理より」と名前を挙げて引用してよいものに限った。
逆に、答案で結論だけ書くと減点されうる箇所には、そのつど注意を書いた。

\section{本書で使う記号}

\begin{definition}{数の集合と論理記号}{tb-notation-sets}
\[
\N=\{1,2,3,\ldots\},\quad
\Z=\{\ldots,-1,0,1,\ldots\},\quad
\Q,\quad \R,\quad \mathbb C
\]
をそれぞれ自然数、整数、有理数、実数、複素数の全体とする。
本書では $0\notin\N$ とし、$0$ を含めたいときは $n\ge0$ と書く。

論理記号は次の意味で使う。
\[
\forall\ \text{（任意の）},\quad
\exists\ \text{（ある〜が存在して）},\quad
\Longrightarrow\ \text{（ならば）},\quad
\Longleftrightarrow\ \text{（同値）}.
\]
$A\setminus B=\{x\in A\mid x\notin B\}$ を差集合、$A\sqcup B$ を
$A\cap B=\varnothing$ が分かっているときの合併（直和）とする。
\end{definition}

\begin{definition}{よく使う記法}{tb-notation-misc}
\[
\abs{x}\ \text{（絶対値）},\quad
\norm{x}\ \text{（ノルム）},\quad
\inner{x}{y}\ \text{（内積）},\quad
\lfloor x\rfloor\ \text{（床関数）}
\]
とする。行列 $A$ の行列式は $|A|$、転置は $A^{\mathsf T}$、
複素共役転置（随伴）は $A^{*}$ と書く。
$I_n$ は $n$ 次単位行列、$O$ は零行列、$\diag(d_1,\ldots,d_n)$ は対角行列である。
写像 $f$ の像を $f(A)$、逆像を $f^{-1}(B)$ と書く。
$f^{-1}$ は逆写像を意味するとは限らないことに注意する。
\end{definition}

\begin{remark}{$\R$ 上の区間}{tb-notation-interval}
\[
(a,b)=\{x\mid a<x<b\},\qquad
[a,b]=\{x\mid a\le x\le b\}
\]
とし、半開区間も同様に書く。$(a,b)$ は座標の組と同じ記号になるが、
文脈で区別する。曖昧になる場面では「開区間 $(a,b)$」と明示する。
\end{remark}

\begin{memo*}
答案で記号を使うときは、問題文で定義されていない記号を断りなく使わないこと。
$\sec$、$\operatorname{adj}$、$\Pi_k$ のような記号は、初出時に一行で定義を書いてから使う。
\end{memo*}
